\section{总体求解框架}

本题的三问并非三个彼此独立的优化问题，而是在同一计算图上逐层加入硬件约束：问题一建立合法拓扑调度与逻辑生命周期模型；问题二保留这些合法性约束，同时允许面向缓存复用重新选择拓扑顺序，并加入真实缓存容量、连续地址分配与必要的 SPILL；问题三进一步引入执行单元、流水并行和地址复用依赖，优化最终执行周期。因此本文采用“逻辑调度—物理内存—时间流水”的统一建模路线。

设计算图为有向无环图 $G=(V,E)$。节点既包括计算/搬运操作，也包括缓冲区的 \texttt{ALLOC}/\texttt{FREE} 管理操作；边 $(u,v)\in E$ 表示节点 $u$ 必须先于 $v$ 执行。缓存类型包括 L1、UB、L0A、L0B、L0C，执行流水包括 CUBE、VECTOR、MTE1、MTE2、MTE3 与 FIXP。全文对节点属性、缓存生命周期和评价指标均使用统一的数据模型，避免三问之间重复建模或出现指标口径漂移。

\begin{figure}[H]
    \centering
    \includegraphics[width=0.95\textwidth]{pipeline_overview.pdf}
    \caption{三问统一求解框架}
    \label{fig:pipeline-overview}
\end{figure}

图~\ref{fig:pipeline-overview} 给出了完整求解链。首先由解析器构造统一 DAG 与 buffer 元数据；随后在问题一中生成多个合法拓扑调度候选，并以峰值驻留量为唯一排序指标；问题二在问题一合法调度空间内引入 Q2-aware 顺序偏好，并对候选顺序执行严格连续地址分配，在无法满足容量与连续区间约束时插入 SPILL；问题三固定问题二已确定的 SPILL 身份、顺序与额外搬运量，通过地址重排、关键流水调度和关键 SPILL 邻域搜索减少执行周期。

本文刻意将“候选生成”和“结果评价”分离。各问的启发式算法可以使用不同的优先级、前瞻或局部搜索，但候选最终必须通过独立 validator/evaluator 重新计算指标后才允许进入正式结果。该设计一方面降低启发式内部估计误差对最终结论的影响，另一方面使不同算法能够在完全一致的题面口径下比较。

从优化层次看，问题一主要对应带前驱约束的列表调度与资源约束排序问题，经典工作已指出优先级规则会显著影响有向无环任务图的调度质量\cite{graham1969multiprocessing,hu1961parallel}。但本文的目标不是直接复现已有列表调度算法，而是结合题目特有的 buffer 生命周期和 L0 单驻留约束构造确定性调度策略，并通过同一 evaluator 对不同策略进行统一筛选。
